\QHeading{解决问题。}
\Problem{① 下面是四（2）班诗朗诵比赛的决赛成绩。}
\par\noindent\begingroup
\begin{tikzpicture}[x=\dimexpr\linewidth-1pt\relax,y=-1mm,draw=black,text=black,line width=.45pt]
\draw (0.000000,0)--(0.000000,24.000);
\draw (0.142857,0)--(0.142857,24.000);
\draw (0.285714,0)--(0.285714,24.000);
\draw (0.428571,0)--(0.428571,24.000);
\draw (0.571429,0)--(0.571429,24.000);
\draw (0.714286,0)--(0.714286,24.000);
\draw (0.857143,0)--(0.857143,24.000);
\draw (1.000000,0)--(1.000000,24.000);
\draw (0,0.000)--(1,0.000);
\draw (0,8.000)--(1,8.000);
\draw (0,16.000)--(1,16.000);
\draw (0,24.000)--(1,24.000);
\node[font=\normalsize,align=center,inner sep=0pt,text width=\dimexpr0.142857\linewidth-2.5mm\relax] at (0.071429,4.000) {姓名};
\node[font=\normalsize,align=center,inner sep=0pt,text width=\dimexpr0.142857\linewidth-2.5mm\relax] at (0.214286,4.000) {$1$号评委};
\node[font=\normalsize,align=center,inner sep=0pt,text width=\dimexpr0.142857\linewidth-2.5mm\relax] at (0.357143,4.000) {$2$号评委};
\node[font=\normalsize,align=center,inner sep=0pt,text width=\dimexpr0.142857\linewidth-2.5mm\relax] at (0.500000,4.000) {$3$号评委};
\node[font=\normalsize,align=center,inner sep=0pt,text width=\dimexpr0.142857\linewidth-2.5mm\relax] at (0.642857,4.000) {$4$号评委};
\node[font=\normalsize,align=center,inner sep=0pt,text width=\dimexpr0.142857\linewidth-2.5mm\relax] at (0.785714,4.000) {$5$号评委};
\node[font=\normalsize,align=center,inner sep=0pt,text width=\dimexpr0.142857\linewidth-2.5mm\relax] at (0.928571,4.000) {最后得分};
\node[font=\normalsize,align=center,inner sep=0pt,text width=\dimexpr0.142857\linewidth-2.5mm\relax] at (0.071429,12.000) {小慧};
\node[font=\normalsize,align=center,inner sep=0pt,text width=\dimexpr0.142857\linewidth-2.5mm\relax] at (0.214286,12.000) {$97$分};
\node[font=\normalsize,align=center,inner sep=0pt,text width=\dimexpr0.142857\linewidth-2.5mm\relax] at (0.357143,12.000) {$98$分};
\node[font=\normalsize,align=center,inner sep=0pt,text width=\dimexpr0.142857\linewidth-2.5mm\relax] at (0.500000,12.000) {$95$分};
\node[font=\normalsize,align=center,inner sep=0pt,text width=\dimexpr0.142857\linewidth-2.5mm\relax] at (0.642857,12.000) {$96$分};
\node[font=\normalsize,align=center,inner sep=0pt,text width=\dimexpr0.142857\linewidth-2.5mm\relax] at (0.785714,12.000) {$99$分};
\node[font=\normalsize,align=center,inner sep=0pt,text width=\dimexpr0.142857\linewidth-2.5mm\relax] at (0.928571,12.000) {{\color{answerred}$97$分}};
\node[font=\normalsize,align=center,inner sep=0pt,text width=\dimexpr0.142857\linewidth-2.5mm\relax] at (0.071429,20.000) {小昆};
\node[font=\normalsize,align=center,inner sep=0pt,text width=\dimexpr0.142857\linewidth-2.5mm\relax] at (0.214286,20.000) {$98$分};
\node[font=\normalsize,align=center,inner sep=0pt,text width=\dimexpr0.142857\linewidth-2.5mm\relax] at (0.357143,20.000) {$94$分};
\node[font=\normalsize,align=center,inner sep=0pt,text width=\dimexpr0.142857\linewidth-2.5mm\relax] at (0.500000,20.000) {$96$分};
\node[font=\normalsize,align=center,inner sep=0pt,text width=\dimexpr0.142857\linewidth-2.5mm\relax] at (0.642857,20.000) {$93$分};
\node[font=\normalsize,align=center,inner sep=0pt,text width=\dimexpr0.142857\linewidth-2.5mm\relax] at (0.785714,20.000) {$100$分};
\node[font=\normalsize,align=center,inner sep=0pt,text width=\dimexpr0.142857\linewidth-2.5mm\relax] at (0.928571,20.000) {{\color{answerred}$96$分}};
\end{tikzpicture}\par\endgroup
\Problem{去掉一个最高分和一个最低分，其他分数的平均数就是最后得分。将上表填写完整，并回答谁是第一名。}
\par\noindent
\Answer{小慧去掉$99$分和$95$分。把$98$分中的$1$分移给$96$分，余下三项均为$97$分。}
\Answer{\Formula{\begin{aligned}[t]&\phantom{={}}(97+98+96)\div3\\
&=(97+97+97)\div3\\
&=97\times3\div3\\
&=97\times(3\div3)\\
&=97\times1\\
&=97\text{（分）}\end{aligned}}}
\Answer{答：小慧的最后得分是$97$分。}



\Answer{小昆去掉$100$分和$93$分。把$98$分中的$2$分移给$94$分，余下三项均为$96$分。}
\Answer{\Formula{\begin{aligned}[t]&\phantom{={}}(98+94+96)\div3\\
&=(96+96+96)\div3\\
&=96\times3\div3\\
&=96\times(3\div3)\\
&=96\times1\\
&=96\text{（分）}\end{aligned}}}
\Answer{答：小昆的最后得分是$96$分。}

\par
\Answer{$97>96$。答：小慧是第一名。}
\par\noindent

\MethodTitle{另一种方法：去掉最高分和最低分后，求总分再平均}

\par\noindent\begin{minipage}[t]{.485\linewidth}
\Answer{小慧：}
\Answer{\Formula{\begin{aligned}[t]&\phantom{={}}(97+98+96)\div3\\
&=(195+96)\div3\\
&=291\div3\\
&=97\text{（分）}\end{aligned}}}
\Answer{答：小慧的最后得分是$97$分。}

\end{minipage}\hfill\begin{minipage}[t]{.485\linewidth}
\Answer{小昆：}
\Answer{\Formula{\begin{aligned}[t]&\phantom{={}}(98+94+96)\div3\\
&=(192+96)\div3\\
&=288\div3\\
&=96\text{（分）}\end{aligned}}}
\Answer{答：小昆的最后得分是$96$分。}

\end{minipage}\par
\Answer{$97>96$。答：小慧是第一名。}

\Problem{② 光明小学三年级有$4$个班，共为地震灾区捐款$720$元；四年级有$3$个班，共为地震灾区捐款$680$元。这两个年级平均每个班捐款多少钱？}
\Answer{\Formula{\begin{aligned}[t]&\phantom{={}}(720+680)\div(4+3)\\
&=1400\div7\\
&=200\text{（元）}\end{aligned}}}
\Answer{答：这两个年级平均每个班捐款$200$元。}




\MethodTitle{另一种核对：用每班200元检验}

\Answer{这里的$200$元是主解已求得的平均数，用下列方法核对，不把猜测当作已知条件。}
\par\noindent\begin{minipage}[t]{.485\linewidth}
\Answer{三年级少的金额：}
\Answer{\Formula{\begin{aligned}[t]&\phantom{={}}200\times4-720\\
&=800-720\\
&=80\text{（元）}\end{aligned}}}

\end{minipage}\hfill\begin{minipage}[t]{.485\linewidth}
\Answer{四年级多的金额：}
\Answer{\Formula{\begin{aligned}[t]&\phantom{={}}680-200\times3\\
&=680-600\\
&=80\text{（元）}\end{aligned}}}

\end{minipage}\par
\Answer{把四年级多出的$80$元补给三年级，两个年级各班都可平均得到$200$元。}
\Answer{答：这两个年级平均每个班捐款$200$元。}

\Problem{③ 方方的一次测试成绩如下：语文$88$分，数学$84$分，英语$80$分。方方平均每科的成绩是多少分？}
\Answer{把语文的$4$分移给英语，三科均为$84$分。}
\Answer{\Formula{\begin{aligned}[t]&\phantom{={}}(88+84+80)\div3\\
&=(84+84+84)\div3\\
&=84\times3\div3\\
&=84\times(3\div3)\\
&=84\times1\\
&=84\text{（分）}\end{aligned}}}
\Answer{答：方方平均每科的成绩是$84$分。}

\par

\MethodTitle{另一种方法：总分除以科数}

\Answer{\Formula{\begin{aligned}[t]&\phantom{={}}(88+84+80)\div3\\
&=(172+80)\div3\\
&=252\div3\\
&=84\text{（分）}\end{aligned}}}
\Answer{答：方方平均每科的成绩是$84$分。}

\Problem{$\star$④ 在一次登山比赛中，亮亮上山时平均每分钟走$30$米，$30$分钟到达山顶；按原路返回时，平均每分钟走$45$米。亮亮往返一次平均每分钟走多少米？}
\Answer{上山路程：$30\times30=900$（米）；下山时间：$900\div45=20$（分）。}
\Answer{\Formula{\begin{aligned}[t]&\phantom{={}}900\times2\div(30+20)\\
&=900\times2\div50\\
&=1800\div50\\
&=36\text{（米）}\end{aligned}}}
\Answer{答：亮亮往返一次平均每分钟走$36$米。}

\MethodTitle{另一种方法：以每分钟30米为基数补足}

\Answer{上山路程$900$米，下山用时$20$分，往返共$50$分。若这$50$分都按每分钟$30$米计算，下山每分钟会少算$15$米。}
\Answer{\Formula{\begin{aligned}[t]&\phantom{={}}30+(45-30)\times20\div50\\
&=30+15\times20\div50\\
&=30+300\div50\\
&=30+6\\
&=36\text{（米）}\end{aligned}}}
\Answer{答：亮亮往返一次平均每分钟走$36$米。}
\Answer{注意：上、下山所用时间不同，不能直接计算$(30+45)\div2$。}
